\scanpage{39}
\uline{\unsure{a'}}
\[
\operatorname{Hom}^d(M^*,\deleted{\unsure{\otimes_k}}N\otimes_k S_k)\simeq\operatorname{Hom}^d(\tau_{N-d}M,\tau_N(N\otimes_k S_k)
\]
\figure{\begin{tikzpicture}\draw (0,0)--(0,1);\draw (.1,1.05)--(8,1.05);\draw (.1,-.05)--(7.7,-.05);\end{tikzpicture}}
si $\left\{\begin{array}{l}M^*\in\mathcal M_{-d}\\N\geq0\end{array}\right.$

p. ex.
\diagram{\begin{tikzcd}[column sep=small]
\operatorname{Hom}^d(M\otimes_k S_k,N\otimes_k S_k) \arrow[r,"\sim"] \arrow[d,phantom,"\wr"] & \operatorname{Hom}^d(\tau_{N-d}M_{S_k},\tau_N N_{S_k} \\
\operatorname{Hom}^{\deleted{\unsure{d}}}_k(M,N) & \text{\uuline{si $d\geq0$}, $N\geq0$}
\end{tikzcd}}
\figure{\begin{tikzpicture}\draw (0,0)--(0,1);\draw (.2,1.1)--(8,1.1);\draw (.3,-.1)--(7.8,-.1);\end{tikzpicture}}

\unsure{Hom. hom} de degré $d\in\mathbb Z$
\diagram{\begin{tikzcd}
\tau_\alpha M_{S_k}\arrow[r] & \tau_\beta N_{S_k}
\end{tikzcd}} $(\beta\geq\alpha)$

a) \deleted{$d\geq0$}. Supposons d'abord $\boxed{d\geq\beta-\alpha}$ i.e $\alpha+d\geq\beta$

On a \deleted{$\unsure{\operatorname{Hom}^i(M)}$} --
\[
\operatorname{Hom}^d(\tau_\alpha M_{S_k},\tau_\beta N_{S_k})\overset{\sim}{\longleftrightarrow}\operatorname{Hom}^d(\tau_\alpha M_{S_k},\underbrace{\tau_{\unsure{\alpha+i|\alpha+d}}\tau_\beta N_{S_k}}_{=\tau_{\alpha+d}N_{S_k}})
\]
\[
\overset{\sim}{\longleftarrow}\operatorname{Hom}^d(M_{S_k},N_{S_k})\simeq\operatorname{Hom}_k(M,N)
\]
(si $u:M-N$, il lui correspond l'hom
\[
\underset{j\geq\alpha}{mT^{(j)}}\longmapsto u(m)T^{(d)}T^{(j)}=c_{d,j}u(m)T^{(d+j)}
\]
)

b) Supposons $0\leq d<\beta-\alpha$ \inserted{$\unsure{\text{i.e.}}\ \beta\geq\beta-d\unsure{>|\geq}\alpha$} on a alors
\deleted{\diagram{\begin{tikzcd}[column sep=small]
\operatorname{Hom}^d(\tau_\alpha M_{S_k},\tau_\beta N_{S_k})\arrow[r,"\unsure{\mathrm{res.}}"'] & \operatorname{Hom}^d(\tau_{\beta-d}M_{S_k},\tau_\beta N_{S_k})\arrow[d,phantom,"\wr"]\\
& \operatorname{Hom}_k(M,N)
\end{tikzcd}}}
\mspage{38}
